\documentclass[11pt,a4paper]{article}

\input{../preambule}

\begin{document}

\title{Semaine 6}
\author{}
\date{09/11 - 15/11 }

\maketitle

\begin{exercise}
    Soit $n$ un carré parfait à quatre chiffres, tous inférieurs ou égaux à 6. Si on ajoute 3 à chacun de ces chiffres, on obtient encore un carré parfait. Trouver $n$.
\end{exercise}

\begin{exercise}
    Trouver tous les nombres premiers $p$ pour lesquels $3p+4$ est un carré parfait
\end{exercise}

\begin{exercise}
    Résoudre l'équation diophantienne~:
        \[x^4 - 2 y^2 = 1\]
    d'inconnue $(x,y) \in \N^2$
\end{exercise}

\begin{exercise}
    Soit $n \in \N$. Montrer que si $n$ est à la fois un carré parfait et un cube parfait, alors il est la puissance sixième d'un entier.
\end{exercise}

\begin{exercise}
\begin{questions}
\item Soit $n\ge 7$. Montrer qu'il existe $a,b\ge 2$ des entiers premiers entre eux tels que $n=a+b$. Discuter de la nécessité de l'hypothèse $n\ge 7$.
\item Soit $k\ge 3$. Notons $p_1<p_2<\ldots$ la suite des nombres premiers. En déduire que $p_{k+1}+p_{k+2}\le p_1 \ldots p_k$.
\item Pour tout $n\ge 2$, on note $q(n)$ le plus petit premier qui ne divise pas $n$. Conclure que $\frac{q(n)}{n}\to 0$.
\end{questions}
\end{exercise}

\begin{exercise}
On note $E=\set{n\ge2\mid n\mid 3^n+1}$
\begin{questions}
\item Montrer que $\forall n\in E, v_2(n)=1$
\item Montrer que $\forall n\in E, 5|n$
\item Montrer que $E$ est infini.
\item Montrer que $5050\in E$
\end{questions}
\end{exercise}

\begin{exercise}
Soit $k\in\N, p\in\mathbb{P}_{\ge4}$. 
Calculer $\displaystyle\sum_{x=0}^{p-1} x^k\pmod{p}$ selon la valeur de $k$.
\end{exercise}

\begin{exercise}
Soit $n\in\N$. Calculer $v_2(5^{2^n}-1)$
\end{exercise}

\begin{exercise}
Soit $A\subset\mathbb{N}$ une partie stable par addition.
Montrer qu'à un nombre fini d'éléments près,
$A=k\mathbb{N}$ pour un certain $k\in\mathbb{N}$.
\end{exercise}

\begin{exercise}
On note $\varphi(n)$ le nombre d'entiers naturels inférieurs
à $n$ et premiers avec $n$.

On admet que $\varphi$ est multiplicative :
\[
    \gcd(n,m)=1
    \implies
    \varphi(nm)=\varphi(n)\varphi(m).
\]
On pose, pour $n\in\mathbb{N}^*$,
\[
    f(n)=\sum_{d\mid n}\varphi(d).
\]
\begin{questions}
    \item Montrer que $f$ est multiplicative.
    \item En déduire que $f=\operatorname{id}_{\mathbb{N}^*}$.
\end{questions}
\end{exercise}

\begin{exercise}
\begin{questions}
    \item Calculer, pour $n\in\mathbb{N}$,
    \[
        v_2\left(5^{2^n}-1\right).
    \]
    \item Soit $p$ un nombre premier impair.
    Comment se comporte la suite
    \[
        \left(v_p\left(5^{2^n}-1\right)\right)_{n\in\mathbb{N}}\,?
    \]
    \item Calculer explicitement cette suite pour $p=3$.
\end{questions}
\end{exercise}

\begin{exercise}
Calculer, pour $n\in\mathbb{N}$,
\[
    10^{10^n}\pmod{7}.
\]
\end{exercise}

\begin{exercise}
Montrer que, pour tous $k,n\in\mathbb{N}$,
\[
    (n!)^k\mid(nk)!.
\]
\end{exercise}

\begin{exercise}
Soit $p$ un nombre premier.
\begin{questions}
    \item Montrer que, pour $1\le k<p$,
    \[
        k\mid\binom{p}{k}.
    \]
    \item Plus généralement, calculer, pour
    $n\in\mathbb{N}^*$ et $1\le k\le p^n$,
    \[
        v_p\left(\binom{p^n}{k}\right).
    \]
\end{questions}
\end{exercise}

\begin{exercise}
On pose, pour $n\in\mathbb{N}^*$,
\[
    f(n)=\prod_{\substack{d\mid n\\d\ne n}}d.
\]
Résoudre l'équation $f(n)=n$.
\end{exercise}
















\end{document}
